\begin{abstract}
Language models are widely used for classification, but generating an answer does not directly
provide a probability distribution over a fixed label set. We evaluate Jev, TypeSafe's first
System One model, zero-shot on four public datasets against zero-shot Qwen3.6-27B self-hosted on
an AWS SageMaker \texttt{\QwenScoringInstance} instance (four \QwenGpu{} GPUs), BERT fine-tuned
on an \texttt{\AwsInstanceTypeBert} instance (one \BertGpu{} GPU), and XGBoost trained on an
\texttt{\AwsInstanceTypeXgb} CPU instance. The comparison measures accuracy, inference cost,
latency, and training-label requirements. On balanced or nearly balanced text tasks evaluated
with equal error costs, Jev reaches \AccPctJevSst\% accuracy on SST-2 and \AccPctJevAgnews\% on
AG News, matching Qwen (\AccPctQwenSst\% and \AccPctQwenAgnews\%) at substantially lower
inference cost and latency. Across the four datasets, Jev costs
\$\CostPerKJevSst--\$\CostPerKJevWdbc{} per 1,000 predictions and takes
\JevLatencyWdbc--\JevLatencySst{} seconds per request, compared with
\$\CostPerKQwenSst--\$\CostPerKQwenWdbc{} in AWS compute and
\QwenLatencySst--\QwenLatencyWdbc{} seconds for Qwen. On AG News, BERT and XGBoost surpass both
zero-shot models with \LabelsToMatchJevBertAgnews{} and \LabelsToMatchJevXgbAgnews{} training
examples, respectively---the smallest tested budgets, alongside validation labels for model
selection. Neither surpasses Jev on SST-2 with all \NTrainSst{} training examples. Once trained,
supervised models offer orders-of-magnitude cheaper inference, with microsecond-to-millisecond
batch time per item. On the imbalanced German Credit and Breast Cancer Wisconsin datasets,
XGBoost achieves the highest default macro-F1 scores, at \MacroFOneXgbGerman{} and
\MacroFOneXgbWdbc, respectively. Jev produces the fewest false negatives but substantially more
false positives. This tradeoff warrants evaluation in applications where missed positive cases
are especially costly, but model selection requires explicit error costs and comparison at
suitable decision thresholds. All Jev results are zero-shot, using one fixed input template per
dataset without prompt optimization. Finally, we evaluate KernelSHAP and LIME, showing that these
standard perturbation-based explanation methods apply directly to Jev's class probabilities and
identify words or features that influence its predictions.
\end{abstract}
